\subsection{窄带随机过程}
\textbf{条件}\quad $\Delta f<<f_c$, $f_c>>0$.

\subsubsection{表示形式}
\textbf{包络相位形式}
\begin{equation}
    \xi(t)=a_\xi(t)\cos(\omega_ct+\varphi_\xi(t)).
\end{equation}

\textbf{同相正交形式}
\begin{equation}
    \xi(t)=\xi_c(t)\cos\omega_ct-\xi_s(t)\sin\omega_ct.
\end{equation}

两者满足三角函数关系:
\begin{gather}
    \begin{cases}
        \xi_c(t)=a_\xi(t)\cos\varphi_\xi(t), \\
        \xi_s(t)=a_\xi(t)\sin\varphi_\xi(t);
    \end{cases} \\
    \begin{cases}
        a_\xi(t)=\sqrt{\xi_c(t)^2+\xi_s(t)^2}, \\
        \varphi_\xi(t)=\arctan\frac{\xi_s(t)}{\xi_c(t)}.
    \end{cases}
\end{gather}

\subsubsection{包络和相位的统计特性}
一个均值为零, 方差为 $\sigma^2$ 的窄带平稳高斯过程 $\xi(t)$, 其包络 $a_\xi(t)$ 的一维分布是瑞利分布, 相位 $\varphi_\xi(t)$ 的一维分布是均匀分布, 且两个分量统计独立.

\subsubsection{同相分量和正交分量的统计特性}
一个均值为零的窄带平稳高斯过程 $\xi(t)$, 其同相分量 $\xi_c(t)$ 和正交分量 $\xi_s(t)$ 同样是平稳高斯过程, 均值为零, 方差也相同, 且两个分量统计独立 (互不相关).
